\documentclass[12pt]{article}
\usepackage[utf8]{inputenc}
\usepackage{amsmath, amssymb, amsfonts}
\usepackage{geometry}
\usepackage{booktabs}
\usepackage{enumitem}
\geometry{margin=1in}
\setlist[enumerate]{leftmargin=*}

\title{Practice Problem Set 2\\ \vspace{15pt} \large ECON 101 -- Macroeconomics}
\author{}
\date{Fall 2026}

\begin{document}
\maketitle


\begin{enumerate}



\item
The market for coffee is initially in equilibrium. Suppose research shows that drinking coffee reduces the risk of heart disease.
\begin{enumerate}
    \item Illustrate on a graph the impact on the demand curve.
    \item Show the new equilibrium price and quantity.
    \item Explain the changes.
\end{enumerate}

\item
In the market for smartphones, a new technology reduces production costs.
\begin{enumerate}
    \item Draw the original supply and demand curves.
    \item Show how the supply curve shifts.
    \item Indicate the effect on equilibrium price and quantity.
\end{enumerate}

\item
Suppose demand and supply in the market for tablets are given by:
\[
Q_d = 100 - 2P \quad \quad Q_s = 20 + 3P
\]
where \(Q\) is quantity and \(P\) is price in dollars.
\begin{enumerate}
    \item Find the equilibrium price and quantity.
    \item Determine whether there is a shortage or surplus if \(P = 10\).
\end{enumerate}

\item
Consider the same market as in Problem 3. Suppose the government imposes a price floor of \(P = 20\).
\begin{enumerate}
    \item Calculate the quantity demanded and supplied at this price.
    \item Determine the size of the surplus.
    \item Is this price floor binding? Why or why not?
\end{enumerate}
\newpage
\item
In the housing rental market, demand and supply are given by:
\[
Q_d = 80 - 2P \quad \quad Q_s = 10 + 2P
\]
If the government sets a rent ceiling at \(P = 15\):
\begin{enumerate}
    \item Find the equilibrium price and quantity without the ceiling.
    \item Compute the shortage created by the ceiling.
    \item Discuss one possible unintended consequence of this policy.
\end{enumerate}




\item \textbf{Expenditure Approach} \\
Suppose in 2024, the following data for Canada are reported (in billions of dollars):  
\begin{itemize}
    \item Consumption = 1{,}400  
    \item Investment = 500  
    \item Government Purchases = 600  
    \item Exports = 400  
    \item Imports = 450  
\end{itemize}

\begin{enumerate}[label=(\alph*)]
    \item Calculate GDP using the expenditure approach.  
    \item If depreciation is \$200 billion and net factor payments from abroad are --\$50 billion, calculate \textbf{Net National Product (NNP)} and \textbf{Gross National Product (GNP)}.  
\end{enumerate}

\item \textbf{Income Approach} \\
A simplified economy reports the following for 2024 (in billions of dollars):  
\begin{itemize}
    \item Wages and Salaries = 2{,}200  
    \item Rental Income = 300  
    \item Corporate Profits = 400  
    \item Net Interest = 200  
    \item Indirect Taxes less Subsidies = 250  
    \item Depreciation = 150  
\end{itemize}

\begin{enumerate}[label=(\alph*)]
    \item Compute GDP using the income approach.  
    \item Compare it to the expenditure method if GDP was \$3{,}350 billion. Are there any statistical discrepancies?  
\end{enumerate}
\newpage
\item \textbf{Real vs Nominal GDP} \\
A two-good economy produces only bread and smartphones.  

\begin{center}
\begin{tabular}{lcccc}
\toprule
Year & Bread (Qty) & Price & Smartphones (Qty) & Price \\
\midrule
2023 & 100 & \$2 & 50 & \$400 \\
2024 & 110 & \$2.20 & 60 & \$380 \\
\bottomrule
\end{tabular}
\end{center}

\begin{enumerate}[label=(\alph*)]
    \item Compute \textbf{Nominal GDP} for 2023 and 2024.  
    \item Using 2023 as the base year, compute \textbf{Real GDP} for 2023 and 2024.  
    \item Calculate the \textbf{GDP deflator} for 2024.  
\end{enumerate}

\end{enumerate}



\begin{enumerate}[resume]

\item \textbf{CPI Calculation} \\
Assume a typical consumer basket contains:  
\begin{itemize}
    \item 10 loaves of bread  
    \item 5 movie tickets  
\end{itemize}

Prices are given below:  

\begin{center}
\begin{tabular}{lcc}
\toprule
Year & Bread Price & Movie Ticket Price \\
\midrule
2023 & \$2 & \$10 \\
2024 & \$2.20 & \$12 \\
\bottomrule
\end{tabular}
\end{center}

\begin{enumerate}[label=(\alph*)]
    \item Compute the \textbf{CPI} in 2023 (base year) and 2024.  
    \item Calculate the \textbf{inflation rate} between 2023 and 2024.  
    \item Compare the CPI inflation rate with the GDP deflator inflation rate from Question 3. Why might they differ?  
\end{enumerate}

\item \textbf{Nominal vs Real Wages} \\
A worker earns \$20 per hour in 2023. The CPI rises from 100 in 2023 to 110 in 2024.  

\begin{enumerate}[label=(\alph*)]
    \item If the worker’s nominal wage in 2024 is \$21, compute the \textbf{real wage} in 2024 (base year = 2023).  
    \item Did the worker experience an increase in purchasing power?  
\end{enumerate}

\end{enumerate}

\end{document}
